\section{Reference Workflow}
\label{sec:workflow}

\subsection{Three trigger points}

The reference workflow runs at three points in the development lifecycle:

\begin{enumerate}
    \item \emph{Per-commit} (push to feature branch). Inference is run on changed files only; verification is not run; the result is informational --- there is no PR comment yet because there is no PR. Cache is updated.
    \item \emph{Per-PR} (PR opened or updated). Inference is run on the diff against the base branch; verification is run on the inferred specs; a markdown summary is posted as a PR comment. Subsequent pushes update the comment in place.
    \item \emph{Per-merge} (merge to main). Whole-codebase re-inference is run; the cache for main is primed for the next round of PRs.
\end{enumerate}

Three triggers, three caches, three decisions. The per-commit cache is a feature-branch-local convenience for the developer; the per-PR cache reuses both the feature branch's per-commit cache and the base branch's per-merge cache; the per-merge cache is the canonical store, populated by the work that has actually been accepted into the trunk.

\subsection{Failure-mode tree}

Inference can fail in several ways. Table~\ref{tab:failure-tree} summarises the reference workflow's response. Critically, \emph{none} of these outcomes blocks the merge by default; all of them surface in the PR comment.

\begin{table}[ht]
\centering
\caption{Failure-mode tree. Every outcome below results in a PR-comment annotation; none blocks the merge by default. Repos that want strict gating opt in via a label-based check.}
\label{tab:failure-tree}
\small
\begin{tabular}{p{4cm}p{8cm}}
\toprule
\textbf{Failure mode} & \textbf{Response} \\
\midrule
Inferrer crashes (engine bug) & Log to PR; do not block. \\
Inferrer times out (>5 min on a single method) & Log to PR; fall back to ``no spec for this method''; do not block. \\
OpenJML rejects an inferred clause & Log clause + assertion to PR; do not block. \\
Inferred spec contains an unsupported pattern & Log warning; emit best-effort clauses; do not block. \\
Successful run & Post a one-line summary to the PR (clauses added / modified / removed; verification pass / fail counts). \\
\bottomrule
\end{tabular}
\end{table}

The optional gating mode is opt-in: a repo that wants a hard gate adds an \texttt{infer-clean-required} label rule to its branch protection. The default ships with no gate.

\subsection{Cache key design}

The cache key for each method is the SHA-256 of:

\begin{lstlisting}[language={}]
sha256(
    method-source-text
  + "|"
  + canonical-form-of-each-callee-spec
  + "|"
  + inferrer-version
)
\end{lstlisting}

The callee-spec hash is what makes the cache compositional-friendly: when a callee's spec changes, every caller's hash changes too, so the cache invalidates precisely the callers without invalidating unrelated methods. The inferrer-version bump invalidates the entire cache on a system upgrade --- the correct default, since a spec inferred by version $N$ may not be expressible by version $N{+}1$.

For per-file caching (a coarser but cheaper variant), the key is the SHA-256 of the file's content plus the inferrer version. This loses the compositional invalidation property but is sufficient when callee analysis is local to the file.

\subsection{What the PR comment looks like}

The summary is collapsed by default so the PR review surface stays clean; expanding shows the per-file outcomes:

\begin{lstlisting}[language={}]
## JML inference summary

**Inference**: 12 inferred, 47 cached, 0 errors.

**Verification**: 12 verified, 0 flagged, 0 errors.

<details><summary>Inferred file outcomes</summary>

```json
{
  "command": "infer",
  "summary": { "inferred": 12, "cached": 47, "errors": 0 },
  "files": [...]
}
```

</details>

<sub>Generated by jml-ci. Specifications are informational; flagged
clauses do not block the merge by default.</sub>
\end{lstlisting}

The format is deliberately minimal. A reviewer who is not interested in the inference output sees one line; one expand-click reveals the detail.
